\subsection{معرفی و انگیزه}
این مکانیزم بر مبنای مبانی نظری پیشرفته در ادبیات توجه و فشرده‌سازی \cite{ref_kmeansquantize_2024} توسعه یافته است.

کوانتایش اسکالر بر پایه خوشه‌بندی، یکی از کارآمدترین روش‌ها برای کاهش حجم حافظه پارامترهای مدل‌های زبانی بزرگ بدون نیاز به دگرگونی در ابعاد فضایی بردارها است. با این وجود، روش‌های مرسوم خوشه‌بندی وزن‌ها با الگوریتم کی-میانگین (\lr{K-Means}) معمولاً ناپیوسته و غیرقابل مشتق‌گیری بوده و اعمال پس‌انتشار خطا برای کالیبراسیون مجدد مقادیر پس از کوانتایش با چالش جدی روبرو است.

ماژول \lr{kmeans\_quantize} با پیاده‌سازی مکانیزم برآوردگر عبور مستقیم (\lr{Straight-Through Estimator - STE}) و ساختار لایه اختصاصی \lr{KMeansLinear} توسعه داده شده است. این ماژول وزن‌های لایه‌های خطی را به $k$ مرکز خوشه تصویر کرده و شاخص‌های حافظه را با نگاشت به صورت شاخص‌های فشرده ذخیره می‌کند. هم‌زمان، به منظور حفظ قابلیت یادگیری و بازیابی دقت، از آداپتورهای با رتبه پایین \lr{LoRA} در کنار مراکز خوشه‌ها بهره می‌گیرد تا بدون تغییر ابعاد شبکه، دقت مدل در سطح بالایی بازیابی شود.
